\section{Consequences for games and markets}\label{sec:weak-nash}

Theorem~\ref{thm:main} supplies the positive failed-gate premise in
the cited reductions to equilibrium problems. We first put its
circuits into the forms those reductions use. This preparation, together
with the finite-witness checks below, makes the transfers precise for
our gate conventions and complete binary encodings. The applications
then concern weak and relative Nash equilibria, allocation and market
equilibria, and stationary equilibria in discounted games. Their
approximation requirements differ, so we specify each target relation
before stating its consequence.

\subsection{Circuit preparation for the applications}

The Boolean tolerance in Definition~\ref{def:gcircuit} is nonmonotone.
We therefore remove Boolean gates before reducing numerical tolerances
or controlling fan-out. The two stages of the following lemma are used
separately below.

\begin{lemma}[Common circuit interfaces]\label{lem:application-circuits}
There are fixed positive rational constants $\eta,\gamma,\tau,\gamma_2$
and deterministic polynomial-time constructions which, from any
End-of-Line instance $I$, produce circuits $G_{\rm nb}$ and $G_2$
using only Constant, Scale, Copy, Add, Subtract, and Less. Their parameters
lie in $[0,1]$, and every used node has a defining gate. The circuit
$G_2$ has fan-out at most two.

Every assignment of $G_{\rm nb}$ satisfying all but a $\gamma$
fraction of its gates at tolerance $\eta$, and every assignment of
$G_2$ satisfying all but a $\gamma_2$ fraction at tolerance $\tau$,
can be converted to an assignment of $R(I)$ meeting the tolerances of
Theorem~\ref{thm:main}. For rational assignments the conversion and the
resulting End-of-Line decoder take time polynomial in the input length
plus the total assignment encoding length.
\end{lemma}

\begin{proof}
Let $\eps_c,\delta_c$ be the constants in Theorem~\ref{thm:main}, and
let $G=(V,T)=R(I)$.

\paragraph{Removing Boolean gates.}
First give every used node a defining gate, if necessary. A node used as
an input but not defined by a gate can be given the self-copy gate
$\mathrm{Copy}(v;v)$, which every assignment satisfies. There are at
most $2|T|$ such nodes, because the original gates have at most $2|T|$
input positions. Thus the new gate count is at most $3|T|$; reducing
the permitted failed-gate fraction by a factor three compensates.
Unused nodes may be omitted and later assigned zero.

On this gate list, use the strong-Boolean relation of Schuldenzucker
and Seuken at the same tolerance $\eps_c$. Their Proposition~2 shows
gatewise that a strong-Boolean $\eps_c$ solution satisfies the
weak-Boolean relations in Definition~\ref{def:gcircuit}, since
$\eps_c<1/10<1/4$. Their Lemma~1 replaces each Boolean gate by a
constant-size gadget with no Boolean gates, such that an
$\eps_c/2$ solution restricts to a strong-Boolean $\eps_c$ solution.
Their Section~6 records the same reduction with a positive failed-gate
fraction: the constant-size gadgets lose only a constant factor in
that fraction~\cite[Proposition~2, Lemma~1, and Section~6]{SchuldenzuckerSeuken2019}.
The comparison gate in that paper has the opposite orientation from
$\mathrm{Less}$ here, so we swap its two inputs. The new parameters
remain bounded rationals.

Consequently there are fixed positive rationals
$\eta=\eps_c/2$ and $\gamma$ and a polynomial-time constructed circuit
$G_{\rm nb}$ using only Constant, Scale, Copy, Add, Subtract, and Less
such that any assignment satisfying its gates at tolerance $\eta$ with
at most a $\gamma$ fraction of failures yields an assignment of $G$
satisfying the original $\eps_c,\delta_c$ bounds. All scale factors
are at most one. This no-Boolean relation is monotone in its numerical
tolerance. We use that monotonicity only after this conversion; at no
point is a weak-Boolean gate reclassified at a larger tolerance.

\paragraph{Fan-out reduction with fixed slack.}
We spell out the precision and failure accounting in the fan-out
steps of~\cite[proof of Theorem~2, pp.~4--5]{BPR2016}.
Suppose $G_{\rm nb}$ has $m$ gates, put
$\gamma_0=\min\{\gamma,1/4\}$, and set
\[
 D=\lceil 8/\gamma_0\rceil,\qquad
 \eta_1=\eta/3,\qquad \gamma_1=\gamma_0/4.
\]
There are at most $2m$ outgoing wire occurrences in total, hence at
most $\gamma_0m/4$ gates with fan-out at least $D$. Delete those gates.
For each wire from a deleted output into a retained gate, supply an
independent $\mathrm{Constant}(0)$ gate. The resulting circuit has
at most $3m$ gates and fan-out less than $D$.

Given an assignment to this circuit, restore each deleted output as
zero. A failed new constant can spoil only its one receiving gate. A
new constant satisfied at tolerance $\eta_1$ changes its receiving
input by at most $\eta_1$ under restoration. For a retained arithmetic
gate satisfied at tolerance $\eta_1$, the residual after changing its
two inputs increases by at most $2\eta_1$; this uses the scale bound
one. For Less, an original input gap greater than $3\eta_1$ remains
greater than $\eta_1$ after those two changes, so its required output
also remains valid at tolerance $3\eta_1=\eta$. Therefore, if $f_1$
gates failed before restoration, at most
$\gamma_0m/4+f_1$ original gates can fail at tolerance $\eta$.
The bound $f_1\leq\gamma_1(3m)$ makes this at most $\gamma_0m$.
This also shows why a smaller numerical tolerance is used when the
replacement constants are made exact. The large-multiplier step in
the cited construction is unnecessary for our bounded scale factors.

Next route the $d>2$ outgoing wires of each remaining gate through
a balanced binary tree of Copy gates. One may use at most $2d$ new
gates and depth at most
\[
 h=\lceil\log_2 D\rceil+1.
\]
If the preceding circuit has $m_1$ gates, there are at most $2m_1$
wire occurrences, so the new fan-out-two circuit has at most $5m_1$
gates. Use the fixed parameters
\[
 \tau=\frac{\eta_1}{2(2h+1)},\qquad
 \gamma_2=\frac{\gamma_1}{5D}.
\]
When every Copy gate on a route is $\tau$-satisfied, its leaf differs
from the root by at most $h\tau$. The same arithmetic and comparison
check shows that a retained gate satisfied at tolerance $\tau$
becomes satisfied at tolerance at most $(2h+1)\tau<\eta_1$ when its
inputs are restored to the roots. A failed Copy gate has fewer than
$D$ descendant receiving gates. Charging these receivers, and charging
each failed retained gate to itself, loses at most $D$ retained gates
per failed gate. Thus at most a $\gamma_2$ fraction of failures in
the new circuit produces at most a $\gamma_1$ fraction in the preceding
one. All depths, blowups, and tolerance losses are fixed independently
of $I$.

Denote the final fan-out-two circuit by $G_2$. All instance maps have
polynomial bit cost: they copy bounded rational parameters and introduce
only fixed rational constants. Their solution maps restrict assignments,
restore zero values, or copy existing coordinates, so they also have
polynomial bit cost on rational inputs.

The resulting assignment of $G$ need not obey the particular witness
polynomial $p(M)$ in Definition~\ref{def:gcircuit}. This causes no
restriction on the decoder. The argument after
\eqref{out:endpoint-proximity} proves soundness for every real assignment
meeting the gate-failure budget, and gives a deterministic decoder for
rational assignments in time polynomial in $N$ plus their total encoding
length. Applying it proves the stated interface.
\end{proof}

\subsection{Weak approximate Nash equilibria}

Babichenko, Papadimitriou, and Rubinstein reduce positive-gap generalized
circuits to weak approximate Nash search, where a fixed fraction of
players may have deviation gain larger than a fixed positive
error~\cite[Definition~4 and Theorem~2]{BPR2016}.
Lemma~\ref{lem:application-circuits} supplies the required circuit
interface. We retain the normalization and finite-witness details of the
transfer.

An $n$-player two-action polymatrix game is specified by an undirected
interaction graph and, for each edge, a pair of rational $2$-by-$2$
payoff matrices. Each player chooses action $0$ or $1$ and receives the
sum of its payoffs on incident edges. We use matrices with entries in
$[0,1/3]$ and graphs of maximum degree at most three, so every total
payoff lies in $[0,1]$. The encoding explicitly lists all players,
edges, and matrices, so its length is at least $n$. For a product
mixed-strategy profile $\sigma$,
write
\[
 \operatorname{reg}_i(\sigma)
 =\max_{a\in\{0,1\}}u_i(a,\sigma_{-i})-u_i(\sigma)
\]
for player $i$'s possible gain from deviation. A rational profile is
encoded by the binary rational probabilities $\sigma_i(1)\in[0,1]$.

\begin{corollary}[Fixed-gap weak approximate Nash hardness]
\label{cor:weak-nash}
There are fixed positive rational constants $\eps_{\rm w},\delta_{\rm w}$
and a fixed polynomial $p_{\rm w}$ for which the following search problem
is PPAD-hard under deterministic polynomial-time many-one search
reductions. Given an explicitly encoded rational two-action polymatrix
game of maximum degree at most three with the normalization above and
complete binary length $L$, find a rational product profile of total
encoding length at most $p_{\rm w}(L)$ such that
\begin{equation}\label{eq:weak-nash}
 \bigl|\{i:\operatorname{reg}_i(\sigma)>\eps_{\rm w}\}\bigr|
 \leq \lfloor\delta_{\rm w}n\rfloor .
\end{equation}
No further regret bound is imposed on the players counted on the left.
\end{corollary}

\begin{proof}
Use the fan-out-two circuit of Lemma~\ref{lem:application-circuits},
with numerical tolerance $\tau$ and failure allowance $\gamma_2$.

\paragraph{The game reduction.}
For a fan-out-two generalized circuit, the gate construction in
Rubinstein~\cite[Section~7]{Rubinstein2018}, used in
the forward reduction of~\cite[Theorem~2]{BPR2016}, has two actions per
player and maximum degree at most three. If the players of a gate
gadget use only actions within $\xi_0$ of a best response, their node
probabilities satisfy that gate at tolerance $2\xi_0$. Choose a fixed
positive rational $\xi_0<\tau/2$. These bounds imply our no-Boolean
gate relations at tolerance $\tau$: an antecedent for Less at
tolerance $\tau$ is also an antecedent at tolerance $2\xi_0$, and
the arithmetic and output bounds are stronger. The scale gadget's
output $\min\{\alpha x,1\}$ equals $\alpha x$ here because $\alpha\leq1$.
Each gate uses a fixed number of players, and a player belongs to only
a fixed number of gate gadgets.
Hence there are fixed integers $C_1,C_2$ such that the game has at most
$C_1$ times as many players as circuit gates and each player not
meeting the well-supported condition can spoil at most $C_2$ gates.
Choose a fixed positive rational
\[
 \delta_{\rm w}\leq\frac{\gamma_2}{C_1C_2}.
\]
The decoded circuit assignment then meets its failure budget whenever
all but a $\delta_{\rm w}$ fraction of players meet that condition.

The gadget matrices have rational entries bounded in absolute value
by a fixed rational $B\geq1$: their variable scale parameters lie in
$[0,1]$, and all other gadget constants are fixed. Replacing every
payoff entry $a$ by $(a+B)/(6B)$ puts the entries in $[0,1/3]$.
For each player this adds an action-independent constant and scales
every deviation gain by $\lambda=1/(6B)$. Put $\xi=\lambda\xi_0$.
Thus the needed well-supported tolerance in the normalized game is
the fixed positive rational $\xi$.

Choose a fixed rational $t>0$ such that
$t(t+13)<\xi$, and put $\eps_{\rm w}=t^2$. We use the local pruning
argument of Rubinstein~\cite[Lemma~5]{Rubinstein2018}; this is the lemma
cited with an earlier number in the proof of BPR's Theorem~2. For every
player whose regret is at most $t^2$, remove actions whose payoff is
more than $(1+1/t)t^2$ below its best response and renormalize.
Leave the other players unchanged. The regret bound ensures that the
removed probability is less than one. The lemma's calculation, with
incoming degree at most three and total payoffs in $[0,1]$, shows that
every retained action of each formerly happy player is within
\[
 t(t+1+4\cdot3)=t(t+13)<\xi
\]
of a best response after all these changes. Unchanged unhappy
neighbors contribute no change to that calculation. The fraction of
players not meeting the well-supported condition therefore does not
increase. This gives the required circuit assignment from any profile
satisfying~\eqref{eq:weak-nash}.

\paragraph{Rational witnesses and decoding.}
All instance maps above have polynomial bit cost. The auxiliary
constants are fixed rationals, and the game matrices copy bounded
rational gate parameters through fixed-size gadgets. For a rational
profile, its regrets require only a fixed number of rational operations
per player. The pruning threshold is rational, so exact comparisons,
deletion of actions, and renormalization have polynomial bit cost.
The extracted node probabilities, restrictions, and restored zero
values consequently form a rational assignment of polynomial total
encoding length.

Lemma~\ref{lem:application-circuits} decodes this assignment in its
total encoding length, so the particular witness polynomial $p(M)$ of
Definition~\ref{def:gcircuit} imposes no additional condition here.

Finally, the finite target search relation has a polynomial witness
bound. Round the probabilities of a mixed Nash equilibrium to a grid
of denominator $q$, where the fixed integer $q$ is greater than
$3/\eps_{\rm w}$. With entries in $[0,1/3]$ and degree at most three,
each pure-action payoff changes by at most $1/q$ when the neighboring
probabilities are rounded. Changing the player's own probability
changes its payoff by at most another $1/q$. Its resulting regret is
therefore at most $3/q<\eps_{\rm w}$. A fixed polynomial $p_{\rm w}$
can encode these $n$ rational probabilities. Satisfaction of
\eqref{eq:weak-nash} is checked by rational arithmetic in polynomial
time. The composed instance map and solution decoder establish the
stated search reduction.
\end{proof}


\subsection{Relative approximate bimatrix equilibria}

For a rational bimatrix game $(A,B)$ with $A\geq0$ and $B\leq0$
entrywise, a pair of mixed strategies $(x,y)$ is a relative
$\eps_{\rm r}$-approximate Nash equilibrium if
\begin{equation}\label{eq:relative-nash}
 x^{\mathsf T}Ay\geq(1-\eps_{\rm r})\max_i(Ay)_i,
 \qquad
 x^{\mathsf T}By\geq(1+\eps_{\rm r})\max_j(x^{\mathsf T}B)_j.
\end{equation}
The plus sign in the second inequality reflects the column player's
nonpositive payoffs: each player's loss is at most $\eps_{\rm r}$
times the absolute value of its best-response payoff. These are
conditions on expected payoffs, not on every action in the support.

\begin{corollary}[Relative approximate Nash hardness]
\label{cor:relative-nash}
There are a fixed rational $0<\eps_{\rm r}<1$ and a fixed polynomial
$p_{\rm r}$ such that the following search problem is PPAD-complete
under deterministic polynomial-time many-one search reductions.
Given an explicitly encoded rational bimatrix game $(A,B)$ with
$A\geq0$, $B\leq0$, and complete binary length $L$, find a rational
mixed-strategy profile of total encoding length at most $p_{\rm r}(L)$
satisfying~\eqref{eq:relative-nash}.
\end{corollary}

\begin{proof}
Apply the implication from Conjecture~2 to Theorem~3 in
\cite[Section~4]{BPR2016}. This published chain passes through a
bipartite degree-three family of weak approximate Nash instances.
Theorem~\ref{thm:main} supplies its fixed-tolerance circuit premise on
a bounded-rational subclass, and Lemma~\ref{lem:application-circuits}
makes the gate-basis and fan-out preparation explicit. No change of
weak-Boolean tolerance is needed. The remaining polynomial reduction
gives hardness for a fixed positive relative error;
choosing a smaller positive rational error only strengthens the
requirements. Its marginal and conditional-probability decoder uses
polynomial-time rational arithmetic. Composing with the circuit
decoder after~\eqref{out:endpoint-proximity}, which is polynomial in
the total rational assignment length, gives the required search
reduction.

For membership, an exact Nash equilibrium of the same game satisfies
\eqref{eq:relative-nash}. Exact rational bimatrix Nash belongs to
PPAD, and every rational bimatrix game has an exact equilibrium of
polynomial encoding length~\cite[Theorem~1.1 and Section~2]{CDT2007Preprint}.
Choose $p_{\rm r}$ to accommodate that bound. The displayed inequalities
are checked by polynomial-time rational arithmetic, so the stated
finite search problem is total and belongs to PPAD.
\end{proof}


\subsection{Indivisible course allocation}

We also apply the course-allocation implication of Babichenko,
Papadimitriou, and Rubinstein~\cite[Theorem~4 and Appendix~B]{BPR2016}.
There are $M\geq1$ courses with nonnegative integer capacities $q_j$
and $N\geq1$ students. Each student's acceptable bundles are explicitly
listed in strict preference order, with the empty bundle as a final,
always-available choice; let $k$ be the maximum bundle size.
At prices $p_j\geq0$ and budgets $b_i>0$, each student receives its
most-preferred bundle $x_i\in\{0,1\}^M$ satisfying $p\cdot x_i\leq b_i$.
Budish's clearing convention and guaranteed threshold are
\cite[Definition~1 and Theorem~1]{Budish2011}
\[
 z_j=\begin{cases}
 \sum_i x_{ij}-q_j,&p_j>0,\\
 \max\{\sum_i x_{ij}-q_j,0\},&p_j=0,
 \end{cases}
 \qquad
 \alpha_* = \frac{\sqrt{M\min\{2k,M\}}}{2}.
\]
For positive budgets, their Gini coefficient is
\[
 \operatorname{Gini}(b)
 =\frac{\sum_{i,h}|b_i-b_h|}{2N\sum_i b_i}.
\]

\begin{corollary}[Low-Gini course-allocation hardness]
\label{cor:course-allocation}
There are a fixed rational $0<\gamma_{\rm c}<1/2$ and a fixed
polynomial $p_{\rm c}$ such that the following search problem is
PPAD-complete under deterministic polynomial-time many-one search
reductions. Given the explicit lists and capacities above, of total
binary length $L$, find rational prices and positive rational budgets,
together with the demanded bundles, of total encoding length at most
$p_{\rm c}(L)$, satisfying
\[
 \norm{z}_2\leq\alpha_*,\qquad
 \operatorname{Gini}(b)\leq\gamma_{\rm c}.
\]
\end{corollary}

\begin{proof}
Use the no-Boolean, fan-out-two circuits of
Lemma~\ref{lem:application-circuits}, with their fixed positive numerical
and failed-gate tolerances. The published reduction in
\cite[Appendix~B.2--B.3]{BPR2016} applies to this bounded-rational
subclass: after common normalization of prices and budgets, any
invalid circuit decoding forces the Gini coefficient to be at least a fixed
positive constant. Choose $\gamma_{\rm c}$ strictly below that
bound. Median normalization and price extraction have polynomial bit
cost on rational outputs. Composing with the circuit decoder after
\eqref{out:endpoint-proximity} therefore gives the hardness reduction,
with time polynomial in the input and total witness length.

For membership, Othman, Papadimitriou, and
Rubinstein~\cite[Appendix~A, Remarks~1--2]{OPR2014} prove that computing
an allocation at this clearing threshold with budgets in
$[1,1+\beta]$ belongs to PPAD in the explicit-list model, for every
fixed rational $\beta>0$. Choose a rational $0<\beta<\gamma_{\rm c}$.
Such budgets have Gini coefficient at most $\beta/2<\gamma_{\rm c}$;
existence is guaranteed by~\cite[Theorem~1]{Budish2011}.
The budget interval selects solutions for membership and is not an
additional requirement of the target problem.

Polynomial-size rational witnesses can also be seen directly.
Fix such an allocation and its zero-price pattern, and cap prices at
$2+\beta$ without changing demands. Keep the budget interval, these
price bounds and affordability of the chosen bundles. Introduce a
common slack $s\in[0,1]$ for the unaffordability of every strictly
preferred listed bundle and for the positivity of every positive
price. Maximizing $s$ over the resulting bounded rational polytope
has a positive optimum and a rational vertex optimum of polynomial
bit length. Demands and the price-sign pattern, hence clearing error,
are preserved. All target conditions are checked by rational
arithmetic and enumeration of the explicit lists; the clearing test
is equivalently $4\sum_jz_j^2\leq M\min\{2k,M\}$.
\end{proof}


\subsection{Markets with approximately optimal bundles}

The market reduction of Deligkas, Fearnley, Hollender, and
Melissourgos~\cite{Fisher2026} uses a ternary circuit intermediary.
We record it first; it will also supply the input to the single-controller
game reduction below.

A Pure-Circuit instance consists of explicit finite lists of nodes and
gates. An assignment $z$ gives each node a value in $\{0,1,\bot\}$.
A gate $\mathrm{NAND}(u,v;w)$ requires $z_w=0$ if $z_u=z_v=1$,
and $z_w=1$ if either input is zero; otherwise it imposes no condition.
A gate $\mathrm{PURIFY}(u;v,w)$ requires at least one of $z_v,z_w$
to lie in $\{0,1\}$, and requires $z_v=z_w=z_u$ whenever
$z_u\in\{0,1\}$.
Each gate uses three distinct nodes, and each node is the output of
exactly one gate. We use the version in which each node is also an input
to exactly one gate~\cite[Section~2.2.1]{Fisher2026}.

\begin{lemma}[The input-once Pure-Circuit interface]
\label{lem:application-pure-circuit}
There are a fixed positive rational $\delta_F$ and a deterministic
polynomial-time reduction from End-of-Line to NAND/PURIFY circuits of
the form just defined. Every ternary assignment satisfying all but a
$\delta_F$ fraction of the gates yields an End-of-Line solution in
polynomial time.
\end{lemma}
\begin{proof}
Apply \cite[Theorem~2.2 and Lemma~A.2]{Fisher2026} to the
no-Boolean, fan-out-two circuits of
Lemma~\ref{lem:application-circuits}. The published bridge gives a fixed positive failure allowance and makes
every node an input to exactly one gate. Its solution decoder returns
rational unary-count values, to which the decoder in
Lemma~\ref{lem:application-circuits} applies. All comparisons and
arithmetic have polynomial bit cost. No change of tolerance for a
weak-Boolean gate is required.
\end{proof}

We now define the simple market families used in the reduction.
There is one unit of each good, and utilities have the separable form
\[
 u_i(x_i)=\sum_j\min\{a_{ij}x_{ij},b_{ij}\},
 \qquad a_{ij}\geq0,\quad b_{ij}\in[0,\infty].
\]
All finite data are rational and explicitly encoded. Each buyer has
nonzero utility for only a bounded number of goods, each good has
nonzero utility for only a bounded number of buyers, and the ratio of
nonzero slopes for each buyer is bounded. These bounds are constants
for the family. Every buyer has at least one uncapped term of positive
slope. Fisher buyers have budget one. In the exchange version, every
buyer initially owns $1/|B|$ units of every good, where $B$ is the set
of buyers; in particular the endowments are identical and the economy
is strongly connected.

For errors $\eps_{\rm m},\delta_{\rm m}>0$, an approximate equilibrium
consists of prices $p_j\geq0$ and bundles $x_{ij}\geq0$. Each bundle
must be affordable and give its buyer at least $1-\delta_{\rm m}$
times the maximum utility attainable within that budget; each good
satisfies $\abs{\sum_i x_{ij}-1}\leq\eps_{\rm m}$. In the exchange
case the budget is the value of the buyer's endowment at $p$.
Affordability is exact. Outputs are rational prices and bundles of
polynomial encoding length; the rational-witness formulation is
discussed in~\cite[Section~2.1]{Fisher2026}.

\begin{corollary}[Approximate market equilibrium hardness]
\label{cor:market-hardness}
For every fixed $0<\eps_{\rm m}<1/9$, there is a fixed
rational $0<\delta_{\rm m}<1$ such that finding an
$(\eps_{\rm m},\delta_{\rm m})$-approximate equilibrium in the simple
Fisher markets above is PPAD-hard under deterministic polynomial-time
many-one search reductions. The same holds for their exchange-market
versions with identical endowments and strongly connected economies.
\end{corollary}

\begin{proof}
Use the circuits of Lemma~\ref{lem:application-pure-circuit}.
For each fixed $\eps_{\rm m}<1/9$,
\cite[Proposition~4.4]{Fisher2026} supplies a fixed positive $\delta_{\rm m}$ and a polynomial-time
reduction to the stated Fisher markets. Composing its solution decoder
with the circuit decoders proves the Fisher assertion of
\cite[Theorem~4.1]{Fisher2026} without its conjectural premise.
For rational $\eps_{\rm m}$, the other constants may be chosen rational
within the reduction's strict margins. For a nonrational fixed
$\eps_{\rm m}$, choose a rational clearing tolerance strictly between
$\eps_{\rm m}$ and $1/9$ and use its reduction: every solution with
the smaller clearing error also satisfies the larger one.
The Pure-Circuit bridge returns rational unary-count values;
for rational market witnesses, all comparisons and decoding have
polynomial bit cost. The decoder after~\eqref{out:endpoint-proximity}
applies in the total length of the resulting rational assignment,
without requiring the particular bound $p(M)$ in
Definition~\ref{def:gcircuit}.

For the exchange markets, normalize prices so that
$\sum_jp_j=|B|$. Every buyer's endowment then has value one, so the
same prices and bundles give a Fisher equilibrium with precisely the
same two errors. The positive-slope uncapped terms ensure that the
price sum is nonzero at any approximate equilibrium. This is the
transfer in~\cite[Appendix~B and Theorem~B.1]{Fisher2026}; normalization
also preserves polynomial rational encoding length.
\end{proof}

\subsection{Stationary equilibria in discounted games}

Daskalakis, Golowich, and Zhang~\cite{DGZ2023} obtain conditional
hardness for stationary equilibria whose incentive constraints are
averaged over an initial distribution. Their premise also follows
from Theorem~\ref{thm:main}.

A finite two-player discounted stochastic game has explicitly encoded
rational transition probabilities, rewards, and initial distribution
$\mu$. A stationary joint policy $\pi$ assigns a distribution on action
pairs to each state, sampled afresh at every visit. Write $\pi_{-i}$
for the other player's marginal policy and use normalized values
\[
 V_i^\pi(\mu)=(1-\gamma)\E_{\mu,\pi}
                 \sum_{t\geq0}\gamma^t r_i(s_t,a_t).
\]
The policy is an $\eps$-stationary Markov coarse correlated equilibrium
(CCE) when
\[
 \sup_{\sigma_i}
 \bigl(V_i^{\sigma_i\times\pi_{-i}}(\mu)-V_i^\pi(\mu)\bigr)
 \leq\eps\qquad(i=1,2),
\]
where $\sigma_i$ ranges over stationary policies chosen independently
of any action recommendation. These constraints use the supplied
$\mu$, rather than requiring the bound at every state. An
$\eps$-stationary Nash equilibrium has the same supplied-$\mu$
incentive bound and additionally requires $\pi$ to be a product policy.

\begin{corollary}[Stationary stochastic-game hardness]
\label{cor:stationary-game-hardness}
For some fixed rational $\eps_{\rm s}>0$, finding an
$\eps_{\rm s}$-stationary Markov CCE is PPAD-hard under deterministic
polynomial-time many-one search reductions, even for two-player
turn-based general-sum games with discount $\gamma=1/2$, rewards in
$[-1,1]$, and at most two actions for the active player at each state.
The inactive player's action affects neither rewards nor transitions.
The same hardness holds for $\eps_{\rm s}$-stationary Nash equilibria
with the same supplied-initial-distribution convention. Outputs are
rational policy tables of polynomial encoding length.
\end{corollary}

\begin{proof}
Use the no-Boolean circuit $G_{\rm nb}$ of
Lemma~\ref{lem:application-circuits}, with its fixed positive rational
tolerance $\eta$ and failure allowance $\delta_{\rm nb}$ (the constant
$\gamma$ of that lemma). The target basis in
\cite[Definition~8 and Appendix~D]{DGZ2023} consists of exact zero/one
constants, clipped linear gates, and comparison gates. For Scale and
Copy, repeat the input with coefficient pairs $(\alpha/2,\alpha/2)$
and $(1/2,1/2)$, respectively; Add and Subtract use $(1,1)$ and
$(1,-1)$. A Constant$(\alpha)$ uses its own exact-one node followed
by the scaled linear gate. Retain Less: each strict source comparison
antecedent implies the corresponding weak target antecedent, and the
target output error is smaller. Choose a positive rational target
tolerance below $\min\{\eta,1/12\}$. Each original gate uses at most
two target gates, and each failed target gate can spoil only its own
original gate. A target failure fraction at most $\delta_{\rm nb}/2$ therefore
suffices. This proves the positive-gap premise in
\cite[Conjecture~10]{DGZ2023}; no weak-Boolean tolerance is changed.

Apply the second assertion of~\cite[Theorem~4]{DGZ2023}. Its turn-based
construction has the stated discount, reward bounds and binary
choices. Taking statewise marginals of a joint policy in such a game
gives a product policy with the same values and unilateral deviation
gains, so its approximate stationary CCE decodes to an approximate
stationary Nash equilibrium with the same error. Choose a smaller
positive rational error if needed.
The fixed reduction parameters can be chosen rational with slack.
For rational policy tables, policy evaluation solves rational linear
systems, and the gate decoder and restrictions have polynomial bit
cost. The decoder after~\eqref{out:endpoint-proximity} then applies in
the total rational assignment length. Polynomial-size rational
approximate policy tables exist by rounding stationary Nash policies
at this fixed discount, as in~\cite[Section~2.5]{DGZ2023}.
\end{proof}

\subsection{Restricted stationary-game families}

The preceding consequence also supplies the input for the triangular
zero-sum construction of Park, Zhang, and Ozdaglar~\cite{PZO2025}.
Here the interaction graph concerns the whole auxiliary function,
not just immediate rewards: for every $V:S\to\mathbb R$, each
\[
 Q_i^V(s,a)=r_i(s,a)+\gamma\sum_{s'}P(s'\mid s,a)V(s')
\]
must be a sum of functions of $(s,a_i,a_j)$ over neighbors $j$ of $i$.
This is the Q-function separability of~\cite[Definition~1]{PZO2025}.
Use the same normalized values and supplied-initial-distribution CCE
inequalities as above, now for all three players, with $\pi_{-i}$ the
joint marginal of the other two players.

\begin{corollary}[Triangular zero-sum stationary-game hardness]
\label{cor:triangle-stationary-hardness}
For some fixed rational $\eps_{\triangle}>0$, finding an
$\eps_{\triangle}$-stationary Markov CCE is PPAD-hard under deterministic
polynomial-time many-one search reductions even for rational three-player
zero-sum games with Q-function-separable interactions on the fixed
triangle $K_3$, discount $1/2$, and each player's stage reward in $[-1,1]$.
The hard games are turn-based, with at most two actions for the active
player and a third player with only one action. Outputs are rational
joint-policy tables of polynomial encoding length; incentives are
required only from the supplied rational initial distribution.
\end{corollary}

\begin{proof}
Apply the triangle construction in~\cite[Appendix~C, Case~1]{PZO2025}
to Corollary~\ref{cor:stationary-game-hardness}.
Write $c(s)\in\{1,2\}$ for the state controller and suppress inactive
actions. Add a one-action third player and retain states, transitions
and initial distribution. Specify oriented edge rewards before scaling by
\[
 (b_{12},b_{13},b_{23})=
 \begin{cases}
 (-r_2,r_1+r_2,0),&c(s)=1,\\
 (r_1,0,r_1+r_2),&c(s)=2,
 \end{cases}
 \qquad b_{ji}=-b_{ij}.
\]
Each term depends only on the actions at its edge. Divide every edge
reward by two. The resulting total rewards are
$(r_1/2,r_2/2,-(r_1+r_2)/2)$, lie in $[-1,1]$, and sum to zero.
For every test function $V$, the continuation term in $Q_i^V$ depends
only on $a_{c(s)}$. Allocate it to edge $\{i,c(s)\}$ when $i\ne c(s)$,
and to any incident edge when $i=c(s)$. This proves the required
separability on $K_3$ for every player and every $V$.

The state process of any new-game policy agrees with that of its
first-two-player restriction in the original game. Each original
player's normalized value and deviation gain in the new game is half
the corresponding quantity in the original game; the dummy has no
deviation. Thus error
$\eps_{\triangle}=\eps_{\rm s}/2$ suffices at the unchanged initial
distribution. The primary source uses unnormalized values, which are
twice our normalized values at discount $1/2$; the stated error
calculation uses our convention throughout. All new data use only
rational addition and halving, and restriction has polynomial bit cost.
Polynomial-size rational witnesses follow by the same fixed-discount
rounding argument as above.
\end{proof}

The single-controller construction of Farina, Kontogiannis, Panageas,
and Pollatos~\cite{FKPP2026} gives a different restricted family.
One player controls transitions at every state, while both players'
actions may affect rewards. Joint policies need not be product policies.

\begin{corollary}[Single-controller stationary-game hardness]
\label{cor:single-controller-hardness}
For some fixed rational $\eps_{\rm sc}>0$, finding an
$\eps_{\rm sc}$-stationary Markov CCE is PPAD-hard under deterministic
polynomial-time many-one search reductions for rational two-player
single-controller games with discount $1/16$, rewards in $[0,1]$,
and at most two actions per player per state. In the hard instances,
the initial law is uniform on the circuit-output variable states of
the cited construction and gives no mass to its absorbing sinks.
Outputs are rational joint-policy tables of polynomial encoding length.
The same assertion, with the same error, holds for every fixed rational
$\gamma\in[1/16,1)$.
\end{corollary}

\begin{proof}
Lemma~\ref{lem:application-pure-circuit} supplies input-once
NAND/PURIFY circuits with a fixed positive failed-gate allowance
$\delta_F$. Convert each NAND$(u,v;w)$ to
\[
 \mathrm{NOT}(u;a),\qquad \mathrm{NOT}(v;b),\qquad
 \mathrm{OR}(a,b;w),
\]
using fresh private nodes $a,b$, and keep PURIFY unchanged.
These are the ternary relations on $\{0,1,\bot\}$
of~\cite[Definition~2.6]{FKPP2026}: if $u=v=1$, the NOT outputs are
both zero and OR forces $w=0$; if either input is zero, its NOT output
is one and OR forces $w=1$, even when the other input is $\bot$.
In all other cases NAND imposes no condition. PURIFY has the same
relation in both bases.
Every node remains the input of exactly one gate, and outputs remain
uniquely defined. An original failed gate charges a failure in its own
disjoint replacement block. Since $m$ original gates give at most $3m$
new gates, a failed-gate allowance $\delta_F/3$ suffices. No numerical
tolerance of a weak-Boolean gate is changed.

This establishes the premise of~\cite[Conjecture~1]{FKPP2026}.
Apply its Theorem~5.5, choosing the fixed reduction parameters and a
positive rational equilibrium error within the strict bounds in its
equations~(23)--(24). Its states consist of circuit-variable states
and constantly many absorbing sinks; the initial law is uniform
precisely on the output-variable states. For a rational joint policy,
equation~(8) of that paper decodes a node from the reward player's
marginal $p$: output zero for $p\leq1/4$, one for $p\geq3/4$, and
$\bot$ otherwise. Marginalization, comparison and restriction have
polynomial bit cost. Follow these by the Fisher unary-count decoder
and the decoder after~\eqref{out:endpoint-proximity}.
At fixed discount, rounding stationary equilibrium policies with
slack gives polynomial-size rational witnesses; policy evaluation and
best-response optimization are rational finite-state discounted MDP
computations. No collapse of correlated policies to Nash policies is
used in this reduction.

Finally, the extension in~\cite[Remark~5.7]{FKPP2026} preserves the
normalized incentive inequalities exactly. Put $\beta=1/16$ and, for
a fixed rational $\gamma\in[\beta,1)$, set
\[
 \tau=\frac{\beta(1-\gamma)}{\gamma(1-\beta)},\qquad
 \widetilde P=(1-\tau)I+\tau P.
\]
For every stationary joint or unilaterally deviated policy $\pi$,
\[
 I-\gamma\widetilde P^\pi
 =\frac{1-\gamma}{1-\beta}(I-\beta P^\pi).
\]
Thus $(1-\gamma)(I-\gamma\widetilde P^\pi)^{-1}r_i^\pi$
equals $(1-\beta)(I-\beta P^\pi)^{-1}r_i^\pi$.
Rewards, actions and initial law are unchanged, and the lazy transition
remains rational and single-controller with polynomial encoding length.
\end{proof}

\subsection{Hylland--Zeckhauser equilibria}

The reduction of Braverman, Liu, Xue, and Zhou~\cite{HZ2026}
also yields constant-error hardness for one-sided allocation with
unit demand. An instance has $n$ agents, $n$ unit-supply goods, and
an explicitly encoded rational utility matrix $u\in[0,1]^{n\times n}$.
In their approximate Hylland--Zeckhauser (HZ) equilibrium
\cite[Definition~2.1]{HZ2026}, an allocation $x\geq0$ and prices
$p\geq0$ satisfy
\[
 \sum_i x_{ij}=1,\qquad \sum_j x_{ij}=1,\qquad \min_jp_j=0,
\]
and, for each agent $i$,
\[
 p\cdot x_i\leq1+\eps,\qquad
 u_i\cdot x_i\geq
 \max\{u_i\cdot y:y\geq0,\ \sum_jy_j=1,\ p\cdot y\leq1\}-\eps.
\]
Thus clearing and unit demand are exact; the additive errors concern
expenditure and utility, with comparison budget one.

\begin{corollary}[Approximate HZ equilibrium hardness]
\label{cor:hz-hardness}
There are a fixed rational $\eps_{\rm h}>0$ and a fixed polynomial
$p_{\rm h}$ such that finding an $\eps_{\rm h}$-HZ equilibrium with
rational prices and allocation of total encoding length at most
$p_{\rm h}(L)$ is PPAD-hard under deterministic polynomial-time
many-one search reductions, where $L$ is the input encoding length.
\end{corollary}

\begin{proof}
Use the fan-out-two circuit of Lemma~\ref{lem:application-circuits},
with its fixed positive numerical tolerance $\tau$ and failure
allowance $\gamma_2$. Its no-Boolean gate relations agree with
\cite[Table~1]{HZ2026}, so no Boolean tolerance is changed.
The reduction underlying~\cite[Theorem~3.2]{HZ2026} has two stages:
Lemma~3.4 of that paper converts these circuits to bounded-degree
threshold games, and Lemma~3.25 converts those games to HZ instances.
The construction uses a fixed integer $m$ and rational market data,
with HZ error $m^{-40}$. In that construction, each threshold-game
vertex $v$ has an associated group $G_{v,1}$ of goods. The first
decoder assigns $v$ the value $\min\{1,m^2p(G_{v,1})\}$, where
$p(G_{v,1})$ is the minimum price in that group. This decoder, the
threshold-game-to-circuit decoder, and the subsequent circuit decoders
have polynomial bit cost on rational witnesses. Composing
with the decoder after~\eqref{out:endpoint-proximity} therefore gives
an End-of-Line solution in time polynomial in the input and total
rational witness length.

For completeness, polynomial-size rational HZ witnesses preserve the
exact equalities above. Start from a normalized exact equilibrium
$(x,p)$, as in~\cite[Lemma~B.2]{HZ2026}; summing budgets and column
sums gives $\sum_jp_j\leq n$.
Round prices upward on a rational grid of mesh at most
$\eps_{\rm h}/4$, retaining zero prices. Express $x$ as a convex
combination of at most $n^2$ permutation matrices and round its
weights, preserving nonnegativity and total weight one, so each row
changes in $\ell_1$ norm by at most $\eps_{\rm h}/[4(n+1)]$.
A denominator of order $n^3/\eps_{\rm h}$ suffices. Row and column
sums remain exact, expenditures increase by less than $\eps_{\rm h}$,
and utilities decrease by less than $\eps_{\rm h}$. Since prices only
increase, the comparison optimum cannot increase. These rational
witnesses have a fixed polynomial encoding bound; their utility
conditions are checked by rational linear programming.
\end{proof}
